% !TEX root = ../main.tex

\bilingual{

}{
  \begin{abstract}

    本文研究有限方差 ($\alpha=2$) 普适性类下,~ 具有一般分支律的 $N$-分支随机游走 ($N$-BRW) 前沿演化的渐近性质. 本文所谓$N$-分支随机游走,~ 是指从$\lfloor N^a\rfloor$ ($a\in[0,1]$) 个粒子出发,~ 每一代根据一般分支律 $\Xi$ (可数无限的点过程) 繁殖,~ 并仅保留所有子代粒子中位置最靠前的 $N$ 个粒子的演化过程.

    主要工作围绕 Brunet--Derrida 关于渐进速度修正项$(\log N)^{-2}$ 标度律的猜想展开：首先,~ 在脊柱分解框架下,~ 利用 Mogul'ski\u{i} 小偏差定理建立水平和倾斜管道的概率估计. ~通过引入最左侧位移的正指数矩假设与分支律交叉项的负指数约束,~ 本文将 A\"id\'ekon 与 Hu 针对特定 Poisson 分支律模型建立的 $L^2$ 轨迹极限定理推广至一般分布. ~证明了在三次对数时间与一次对数空间缩放尺度下,~ 前沿位置收敛于线性轨迹 $-a + \frac{\pi^2}{2}\sigma^2 t$,~ 确立了该轨迹作为系统本征属性的普适性. ~其次,~ 针对单粒子起始情形 ($a=0$) 改进了证明策略,~ 利用超临界性质通过位移截断构造至多二分叉过程,~ 并引入``重启机制''克服截断后的灭绝风险,~ 保障了极限轨迹在主尺度下的稳定性.

    最后,~ 利用上述极限定理,~ 本文严格验证了一般分支律下速度修正项的 $(\log N)^{-2}$ 标度律,~ 并给出了包含显式常数的渐近速度表达式. ~本文结果精细刻画了前沿演化的动态过程,~ 为理解有限种群选择机制下的微观涨落与宏观速度偏移之间的关联提供了完备的数学基础.
  \end{abstract}
}

\ifchinese
  \begin{abstract*}
\else
  \begin{abstract}
\fi
  This thesis investigates the asymptotic front evolution of the $N$-branching random walk ($N$-BRW) within the finite variance ($\alpha=2$) universality class with general branching laws. The $N$-branching random walk considered here starts with $\lfloor N^a\rfloor$ ($a\in[0,1]$) particles, where each generation reproduces according to a general branching law $\Xi$ (a countably infinite point process), and only the $N$ leftmost particles among all offspring are retained for the subsequent evolution.

  This thesis mainly focuses on the Brunet--Derrida conjecture regarding the speed correction: First, within the framework of spine decomposition, the Mogul'ski\u{i} small deviation theorem is utilized to establish small deviation estimates in horizontal and tilted tubes. By imposing a positive exponential moment assumption on the leftmost displacement and negative exponential constraints on the branching cross-terms, the $L^2$ trajectory limit theorem of A\"id\'ekon and Hu has been generalized to general branching laws. It has been proven that under triple-logarithmic time and single-logarithmic space scaling, the front position converges to a linear trajectory $-a + \frac{\pi^2}{2}\sigma^2 t$, confirming its universality as an intrinsic property of the system. Second, for the case of a single-particle start ($a=0$), the proof strategy is enhanced by constructing an at-most binary branching process via displacement truncation and implementing a "re-booting mechanism" to mitigate extinction risks, ensuring the stability of the limit trajectory under the primary scale.

  Finally, based on the aforementioned limit theorems, it has been rigorously verified that the $(\log N)^{-2}$ scaling law of the speed correction under general laws holds and the explicit constant is derived. These results provide a refined characterization of front dynamics and offer a comprehensive mathematical foundation for understanding selection mechanisms in finite populations.
\ifchinese
  \end{abstract*}
\else
  \end{abstract}
\fi